\cotefile{157-5}
\scanpage{117}
\mspage{104}
Il reste à prouver (6.1.12), et on procède
comme pour 3.3 (p.~54), en utilisant
le fait qu'une flèche $f$ dans $M(I)_c$
qui est dans $W(I)_{\deleted{\unsure{c}}}$, s'écrit \inserted{dans $M(I)_c$} sous la
forme $f=pi$, où $i\in TC(I)$, et
où $p$ est \deleted{\unsure{l'inv.}} inverse à gauche (rétraction)
d'une flèche $j\in TC(I)$. On aura
donc $u_!(f)=u_!(p)u_!(i)$, \deleted{\unsure{et}} $u_!(i)\in TC(I')$
par (6.\unsure{1}.10), et $u_!(p)$ est inverse à
gauche de $u_!(j)\deleted{\unsure{\in}} \in TC(I')$. On a
donc $u_!(i)$, $u_!(j)$, $u_!(p)\in W(I')$, d'où
$u_!(f)\in W(I')$, qed.

Prouvons enfin d). \deleted{\unsure{Il}} Compte
tenu de la définition de $C_I$, pour
établir $C(I)\subset C_I$, il \deleted{\unsure{est} \illegible{1}} \inserted{suffit de}
voir ceci

% slanted
\margin{\uline{NB} Ce résultat se généralise sous les \unsure{condit\textsuperscript{ns}|conditions} \unsure{inscrites|précisées} dans 6.1.5.}
\uline{Corollaire} 6.2. Soit $I$ \unsure{quelc.} dans \unsure{Cat},
soient $U\subset V$ deux fermés dans $I$
(i.e. des ouverts dans $I^\circ$), et \deleted{\unsure{considé}}
soit $\deleted{f}\unsure{u|f}\colon X_\cdot\longrightarrow Y_\cdot$ une flèche
dans $\deleted{\unsure{f}}C(I)$ (resp. dans $TC(I)$).
\deleted{Considérons la \unsure{fl}} Posons alors
% archive sheet number: 116

\scanpage{118}
\mspage{105}
\[
 X_U=\lim_{i\in U}X(i),\qquad X_V=\underset{i\in V}{\unsure{\underrightarrow{\lim}|\underline{\lim}}}\,X(i),
\]
et de $=$ pour $Y$, et considérons le
diagramme

\diagram{\begin{tikzcd}[ampersand replacement=\&,column sep=large]
\& X_U \arrow[r,"\text{trans.}"] \arrow[d,"f_U"'] \arrow[dr,phantom,"\text{cocart}"] \& X_V \arrow[d] \arrow[dd,bend left=48,dashed,"f_V"]\\
(6.2.1) \& Y_U \arrow[r,"{\deleted{\unsure{\text{trans.}}}}"'] \arrow[dr,dashed,"\text{trans.}"'] \& Y_{U,V} \arrow[d,"{f_{U,V}}"']\\
\& \& Y_V
\end{tikzcd}}

Alors la flèche $u_{U,V}$ est dans $C$ (resp. $TC$).

Démonstration de 6.2. (\uuline{Admettons} \deleted{$U,V$}
que la restriction d'une flèche de
$\deleted{\unsure{f}} C(I)$ resp $TC(I)$ à un fermé
\deleted{\unsure{de}} $\deleted{\unsure{U}}V$ de $I$ \deleted{de} est dans $\deleted{\unsure{f}} C(V)$ \unsure{respt.}
$TC(V)$ Cf. \ldots\ ci-dessous). Cela
nous permet de supposer $V=I$.

La donnée de $U$, qui est un ouvert
de $\unsure{I|I^\circ}$, équivaut à la donnée d'une
flèche \unsure{$v$|``} $I^\circ\longrightarrow\Delta^1$ (\deleted{\unsure{par}} prenant la
fibre en $0$), ou encore à la donnée
d'une flèche $u\colon I\unsure{\longrightarrow|\text{---}}(\Delta^1)^\circ\simeq\Delta^1$
(en prenant la fibre en $1$), \unsure{soit}
\[
 u\colon I\longrightarrow\Delta^1,\qquad
 U=\deleted{\unsure{V}}I_1\overset{\mathrm{def}}{=}u^{-1}(\{1\}).\tag*{(6.2.2)}
\]
% archive sheet number: 117

\scanpage{119}
\mspage{106}
On a alors
\[
 X_U=u_!(X)(1),\qquad X_I=u_!(X)(0)
\]
et le morphisme de transition
$X_U\longrightarrow X_I$ est déduit de \inserted{la flèche} $0\xrightarrow[\unsure{\delta|d}]{}1$ dans
$\Delta^1$, et de $=$ pour $Y$. Soit
\[
 \unsure{\xi|\zeta}=u_!(X),\qquad \eta=u_!(Y)
\]
donc le diagramme (6.2.1) s'identifie
au diagramme déduit de
\[
 \varphi\colon \xi\longrightarrow\eta,\qquad \varphi=u_!(f),
\]
comme

\diagram{\begin{tikzcd}[ampersand replacement=\&,column sep=large]
\& \xi(1) \arrow[r,"{\xi(\delta)}"] \arrow[d,"{\varphi(1)}"'] \arrow[dr,phantom,"{\unsure{\text{cocart}|\text{cart}}}"] \& \xi(0) \arrow[d,"{\deleted{\unsure{\varphi(0)}}}"] \arrow[dd,bend left=50,preaction={draw,dash pattern=on .5pt off 6pt,line width=1.1pt,-},"{\varphi(0)}"]\\
(6.2.3) \& \eta(1) \arrow[r,"{\deleted{\unsure{\eta(\delta)}}}"] \arrow[dr,dashed,"{\eta(\delta)}"'] \& \widetilde{\eta}_{\unsure{1}}\deleted{\unsure{\eta_{1,0}}} \arrow[d,"{\widetilde{\varphi}\deleted{\unsure{\varphi_{1,0}}}}"']\\
\& \& \unsure{i}\eta(0)
\end{tikzcd}}

Or par 6.1.10 appliqué à $u\colon I\unsure{\text{---}|\longrightarrow}\Delta^1$,
la \unsure{relation|condition} $f\in C(I)$ resp $f\in TC(I)$
implique $\varphi\in C(\Delta^1)$ resp. $\varphi\in TC(\Delta^1)$.
Donc pour établir 6.1, on est ramené
au cas particulier où $I=\Delta^1$, $U=\{1\}$, $V=I$,
donc à prouver le

\uline{Lemme 6.2.4.} Soit $\varphi\colon\xi\longrightarrow\eta$ une
flèche dans $M(\unsure{\Delta^1|\Delta})$. Si $\varphi\in C(\Delta^1)$ (resp.
% archive sheet number: 118

\scanpage{120}
\mspage{107}
$\varphi\in TC(\Delta^1)$) alors la flèche $\deleted{\unsure{f_{1,0}}}\widetilde{\varphi}$ dans
le diagramme (6.2.3) est dans $C$ (resp. $TC$).

Plus généralement et plus précisément :

\uline{Lemme} 6.2.5. Soit $M$ une catégorie \inserted{\unsure{stable} par sommes \unsure{amalg.|ind.}},
et $\Phi$ une partie $Q$-close à gauche
(on prendra $\Phi=C$, et $\Phi=TC$, dans le cas
du lemme 6.2.4). \deleted{\unsure{Munissons}} Soit
$\Psi=\Phi_*$, et considérons, pour \unsure{tte} \deleted{\unsure{catég}}
$I$ dans $\mathrm{Cat}$, $\Phi(I)$, $\Psi\langle I\rangle$ dans $M(I)$
\unsure{définis|définies} par
\[
 \begin{aligned}
 &u\colon X\longrightarrow Y\ \inserted{\text{dans }M(I)}\ \text{est dans }\Psi\langle I\rangle\ \text{sss}\\
 &\text{\unsure{ssi|sss}}\quad \forall i\in\operatorname{Ob} I,\quad u(i)\colon X(i)\unsure{\text{---}|\longrightarrow}Y(i)\\
 &\text{est dans }I,
 \end{aligned}\tag*{(6.2.5.1)}
\]
\[
 \Phi(I)=\Psi\langle I\rangle^*.\tag*{(6.2.5.2)}
\]
Prenons le cas $I=\Delta^1\inserted{=\{0\xrightarrow{\unsure{d|\delta}}1\}}$, et soit
$\deleted{\unsure{f}}u\colon X\longrightarrow Y$ une flèche dans $M(\Delta^1)$
correspondant à un carré comm.

\diagram{\begin{tikzcd}[ampersand replacement=\&]
\& X(1)\arrow[r,"{X(d)}"]\arrow[d,"{u(1)}"'] \& X(0)\arrow[d,"{u(0)}"]\\
(6.2.5.3)\& Y(1)\arrow[r,dash,"{Y(d)}"] \& Y(0).
\end{tikzcd}}

Considérons le \unsure{diagramme|codiagramme} correspondant
% archive sheet number: 119

\scanpage{121}
\mspage{108}
\diagram{\begin{tikzcd}[ampersand replacement=\&,column sep=large]
(6.2.5.4) \& X(0)\arrow[r,dash,"{X(\unsure{d|\alpha})}"]\arrow[d,"{u(0)}"']\arrow[dr,phantom,"\text{cocart}"] \& \unsure{x|X}(1)\arrow[d]\arrow[dd,bend left=48,dash,dashed,"{u(1)}"]\\
\& Y(0)\arrow[r,"{\unsure{v|Y}}"]\arrow[dr,dashed,"{Y(\alpha)}"'] \& \widetilde Y\arrow[d,"\widetilde u"']\\
\& \& Y(1)
\end{tikzcd}}

Pour qu'on ait $u\in\Phi(\Delta^1)$, il f et s
que l'on ait
\[
 u(0),u(1),\widetilde u\ \in\ \Phi.\tag*{(6.2.5.5)}
\]

J'ai déjà vérifié la \unsure{chose} (sous forme
duale), mais ne retrouve plus l'endroit
(peut-être n'ai-je pas pris des notes).
Je n'ai pas envie de rechercher la
démonstration ici.

Pour achever de prouver 6.1. d),
i.e. l'égalité $C(I)=C_I$ pour $I$
artinien, j'ai envie de prouver encore
\unsure{qque} chose de plus général que voici :

\uline{Corollaire} 6.3. Soient $(M,\Phi,\Psi)$ \inserted{et $I\longmapsto\Phi(I)\unsure{\subset|\in}\operatorname{Fl}M(I)$} comme
dans 6.2.5. Soit $I$ un ens. ordonné,
\deleted{Alors \unsure{on}} \inserted{et} considérons aussi $\Phi_I$ défini (dans la
\deleted{(6.2.\unsure{1|2}} situation duale) \deleted{\unsure{page}} 25--(2.2) (mais
avec des notations duales, \unsure{i.e.}). Je dis que l'on
\[
 \Phi_I\supset\Phi(I),\tag*{(6.3.1)}
\]
% archive sheet number: 120

\scanpage{122}
\mspage{109}
et \deleted{\unsure{on}} a égalité si $I$ est \unsure{noethérien}, i.e. $I^\circ$ est
artinien.

\uline{Dém}. La démonstration de $\deleted{\unsure{\Phi_I}}\Phi(I)\subset\Phi_I$
est celle qu'on vient de donner. Il reste
à prouver que si $I$ est \deleted{artinien}\inserted{noethérien}, on a
$\Phi_I\subset\Phi(I)$. Etant donné $\unsure{i|u}\colon A\longrightarrow B\unsure{=}$
dans $\Phi_I$, il faut prouver que $u\in\Phi(I)$, i.e.
que pour un diagramme comm. dans $M(I)$,

\diagram{\((6.3.2)\)\quad
\deleted{\begin{tikzcd}[ampersand replacement=\&]
A\arrow[r,"i"]\arrow[d] \& B\arrow[d]\arrow[dl,dashed]\\
X\arrow[r,"p"] \& Y
\end{tikzcd}}\qquad
\begin{tikzcd}[ampersand replacement=\&]
A\arrow[r,"\alpha"]\arrow[d,"i"'] \& X\arrow[d,"p"]\\
B\arrow[r,"\beta"']\arrow[ur,dashed,"h"] \& Y
\end{tikzcd}}

avec $p\in\Psi\langle I\rangle$ i.e. $p(i)\colon X(i)\unsure{\text{---}|\longrightarrow}Y(i)$
dans $\Psi$ pour \unsure{tt} $i$, $\exists h\colon B\longrightarrow X$ laissant le
diagramme commutatif. On va revenir
aux notations duales, en travaillant dans
$M(I^\circ)$, avec $I$ artinien. On considère
l'\unsure{ens} \inserted{$\unsure{\mathcal H|H}$} des couples $(U,h_U)$ d'un ouvert $U$
de $I$ et d'un $h_U\colon B|U\longrightarrow X|U$,
rendant commutatif la restriction à $U$
du \unsure{car.|carr.} On ordonne $\mathcal H$ par prolongement,
on voit que tte \deleted{partie} tot. ordonnée \inserted{$(U_\alpha,h_{U_\alpha})$} dans
$\mathcal H$ a une borne supérieure, savoir
$U=\bigcup_\alpha U_\alpha$, $h_U$ défini \deleted{par les} $h_{U_\alpha}$ sur les $U_\alpha$.
% archive sheet number: 121

\scanpage{123}
\mspage{110.}
\inserted{(N\uline{B}.} \deleted{\unsure{Par Zorn.}} $\unsure{\mathcal H|H}$ n'est pas vide, car $(\varnothing,h_\varnothing)$ est
dans $\mathcal H$. -- mais c'est contenu dans l'\deleted{\unsure{énoncé}}
énoncé précédent que tte partie \unsure{tot} ordonnée
de $\mathcal H$ y a une borne sup\ldots) Par Zorn
il existe un élément maximal $(U,h_U)$.
À prouver que $U=I$, et on aura terminé.
\inserted{On \unsure{va voir|verra} que si $(U,h_U)$ tel que $U\ne I$, ce n'est \uline{pas} maximal.}
Soit $\deleted{\unsure{B}}F=I\setminus U\ne\varnothing$ le fermé complémentaire.
Par l'hypothèse artinienne sur $I$,
$F$ a un élément minimal $i_0$. \deleted{On}
Soit
\[
 U'=U\cup\overbrace{I_{/\unsure{j_0|i_0}}}^{V}=U\cup\{i_0\},
\]
\deleted{on} \deleted{\unsure{doit}}\inserted{\unsure{si}} prolonger $h_U$ en $h_{U'}$ sur $U'$,
on aura terminé. Pour ceci, soit
\deleted{on est ramené à faire la}
$V_0=V\cap U=V\setminus\{j_0\}$, on est ramené
à prolonger la restriction \inserted{$h_{V_0}$} de $h_U$ à $V_0$,
de $V_0$ à $V$. Au total, \unsure{OPS}
$I=I_{/i_0}$ muni d'un élément final $i_0$,
$U=I\setminus\{i_0\}$, et montrer qu'un
$(U,h_U)\in\mathcal H$ se prolonge à $I$. Soient
$A',B',X',Y'$ les restrictions de $A,B,X,Y$
à $U$ \deleted{\unsure{et}} (qu'on pourra aussi noter $I'$).
La donnée de $A$ équivaut à celle
de $A'$, celle de $A_0=A(i_0)\in\operatorname{Ob} M$, et
d'une flèche de
\[
 A'_{I'}\overset{\mathrm{def}}{=}
 \underset{i'\in I'}{\underrightarrow{\lim}}\,A'(i')\xrightarrow{\unsure{\lambda_A|\partial_A}}
 A_0\bigl(\simeq\underset{\unsure{\nearrow}}{A_I}\bigr),
\]
% archive sheet number: 122

\scanpage{124}
\mspage{111}
et de $=$ pour $B,X,Y$. Plus précisément,
la catégorie $M(I^\circ)$ s'interprète comme
celle des triples $X=(X',X_0,\alpha_X)$, où
$X'\in M(I'{}^\circ)$, $X_0\in M$, et $\alpha_X\colon X'_{I'}\longrightarrow X_0$.

À ce compte, la donnée du diagramme
(6.3.2) équivaut à celle d'un diagramme \inserted{carré}
\unsure{similaire|semblable}

\diagram{\((6.3.3)\)\qquad
\begin{tikzcd}[ampersand replacement=\&]
 A'\arrow[r,"{\alpha'}"]\arrow[d,"{i'}"'] \& X'\arrow[d,"{p'}"]\\
 B'\arrow[r,"{\beta'}"]\arrow[ur,dotted,"{h'}"] \& Y'
\end{tikzcd}}

\uline{dans $M(I'{}^\circ)$}, \unsure{d'où|d'un} \unsure{autre|un} diagramme
dans $M$

\diagram{\((6.3.4)\)\qquad
\begin{tikzcd}[ampersand replacement=\&]
 A'_{I'}\arrow[r,"{\alpha'_{I'}}"]\arrow[d,"{i'_{I'}}"'] \& X'_{I'}\arrow[d,"{p'_{I'}}"]\\
 B'_{I'}\arrow[r,"{\beta'_{I'}}"]\arrow[ur,dotted,"{h'_{I'}}"] \& Y'_{I'},
\end{tikzcd}}

d'un \inserted{autre} diagramme carré \uline{dans $M$}

\diagram{\((6.3.5)\)\qquad
\begin{tikzcd}[ampersand replacement=\&]
 A_0\arrow[r,"{\alpha_0}"]\arrow[d,"{i_0}"'] \& X_0\arrow[d,"{p_0}"]\\
 B_0\arrow[r,"{\beta_0}"] \& Y_0
\end{tikzcd}}

et enfin d'un hom. \inserted{$\lambda=(\lambda_A,\lambda_B,\lambda_X,\lambda_Y))$} du carré (6.3.4)
dans le carré (6.3.5.) : \deleted{\unsure{la}}
% archive sheet number: 123

\scanpage{125}
\mspage{112}
\diagram{\((6.3.6)\)\quad
\begin{tikzcd}[ampersand replacement=\&,column sep=small,row sep=small]
 A'_{I'}\arrow[rr,"{\alpha_{I'}}"]\arrow[dr,"{i'_{I'}}"']\arrow[dd,"{\lambda_A}"'] \& \& X'_{I'}\arrow[dr,"{p_{I'}}"]\arrow[dd,"{\lambda_X}"{pos=.8}] \& \\
 \& B'_{I'}\arrow[rr,"{\beta_{I'}}",crossing over]\arrow[ur,dashed,"{h'_{I'}}"]\arrow[dd,"{\lambda_B}"{pos=.75},crossing over] \& \& Y'_{I'}\arrow[dd,dash,"{\lambda_Y}"]\\
 A_0\arrow[rr,"{\deleted{\unsure{\alpha}}}"{pos=.2},"{\alpha_0}"{pos=.75}]\arrow[dr,"{i_0}"'] \& \& X_0\arrow[dr,"{p_0}"] \& \\
 \& B_0\arrow[rr,"{\beta_0}"]\arrow[ur,dotted,"{h_0}"] \& \& Y_0
\end{tikzcd}}

\margin{NB. La condition $p\in\Psi\langle I\rangle$ s'exprime par $p'\in\Psi(I')$ et $\boxed{p_0\in\Psi}$}
D'autre part, la question de prolonger
$(I',h')\in\unsure{\mathcal H|H}$ en $(I,h)$, revient à celle
de trouver $h_0\colon B_0\longrightarrow X_0$ dans $M$,
de façon qu'avec $h'_{I'}$ \unsure{cel|cela} fasse encore
un diagramme commutatif \unsure{de|dans} (6.3.6),
ce qui revient à dire que \inserted{d'une part} les deux
triangles formés dans le carré du
bas sont commutatifs :
\[
 p_0 h_0=\beta_0,\qquad h_0 i_0=\alpha_0,\tag*{(6.3.7)}
\]
et d'autre part le carré diagonal défini
par $h'_{I'}$, $h_0$ soit commutatif (ce qui
exprime qu'on a une flèche $h$ de
$B$ dans $X$) :
\[
 \lambda_X h'_{I'}=h_0\lambda_B.\tag*{(6.3.8)}
\]
% archive sheet number: 124

\scanpage{126}
\mspage{113}
Mais le diagramme (6.3.6), où la flèche
en traits discontinus $h'_{I'}$ est déjà acquise,
et $h_0$ en pointillés reste à trouver, donne
naissance au diagramme, dans lequel
est \unsure{mis|mise} en évidence la flèche \uline{$\widetilde i\in\Phi$}
qui exprime l'hypothèse $i\in C(I)$,
compte tenu de $i'\in C(I')$ :

\diagram{\((6.3.9)\)\quad
\begin{tikzcd}[ampersand replacement=\&,column sep=small,row sep=small]
 A'_{I'}\arrow[rr,"{\alpha_{I'}}"]\arrow[dr,"{i'_{I'}}"]\arrow[dd,"{\lambda_A}"'] \& \& X'_{I'}\arrow[dr,"{p_{I'}}"]\arrow[dd,"{\lambda_X}"{pos=.23}] \& \\
 \& B'_{I'}\arrow[rr,"{\beta_{I'}}",crossing over]\arrow[ur,dashed,"{h'_{I'}}"]\arrow[dd,"s"'{pos=.77},crossing over] \& \& Y'_{I'}\arrow[dd,"{\lambda_Y}"]\\
 A_0\arrow[rr,"{\alpha_0}"{pos=.78}]\arrow[dr,"r"']\arrow[ddr,dashed,"{i_0}"']\arrow[ur,phantom,"{\text{\uline{cocart}}}"{pos=.35}] \& \& X_0\arrow[dr,"{p_0\in\Psi}"] \& \\
 \& \widetilde B\arrow[rr,preaction={draw,solid,line width=.35pt,-},dashed,very thick,"{\widetilde\beta_0}"]\arrow[ur,dashed,thick,"{\widetilde h_0}"]\arrow[d,"{\deleted{\unsure{i_0}}\widetilde i\in\Phi}"'] \& \& Y_0\\
 \& B_0\arrow[uur,dotted,"{h_0}"]\arrow[urr,"{\beta_0}"'] \& \&
\end{tikzcd}}

où la flèche $\widetilde h_0$ en traits discontinus,
de $\widetilde B\overset{\mathrm{def}}{=}A_0\amalg_{A'_{I'}}B'_{I'}$ dans $X_0$, est définie
par les conditions de commutativité
\[
 \widetilde h_0r=\alpha_0,\qquad \unsure{\bar h|\widetilde h}_0s=\lambda_X h'_{I'}\tag*{(6.3.10)}
\]
ce qui est licite, vu que
\[
 \begin{aligned}
 \alpha_0\lambda_A &=
 \underset{\Vert}{\underline{(\lambda_X h'_{I'})i'_{I'}}}\qquad\deleted{\unsure{=}}\\[-2pt]
 &\phantom{=}\underbrace{\lambda_X\bigl(\underset{\alpha_{I'}}{\underline{h'_{I'}i'_{I'}}}\bigr)}_{\substack{\Vert\\\lambda_X\alpha_{I'}}}
 \end{aligned}
\]
% archive sheet number: 125

\scanpage{127}
\mspage{114}
Quant à la flèche $\widetilde\beta_0\colon\widetilde B\unsure{\text{---}|\longrightarrow}Y_0$, elle est
donnée par les conditions de commutativité
\[
 \widetilde\beta_0r=\unsure{p_0|\beta_0}\alpha_0,\qquad
 \widetilde\beta_0s=\lambda_Y\beta_{I'},\tag*{(6.3.11)}
\]
ce qui est \inserted{encore} licite, puisque
\[
 (p_0\alpha_0)\lambda_A=(\lambda_Y\beta_{I'})i'_{I'}
\]
en vertu de la commutativité du
cube (6.3.6). Pour établir la commutativité
totale du diagramme (6.3.9)\unsure{,}
(traits pleins, et \deleted{traits}\unsure{,} discontinus
pour $h'_{I'}$ \deleted{et}, $\widetilde h_0$, et $\widetilde\beta_0$ -- à l'exclusion
de $h_0$ pas \unsure{encore|encor} construit), il faut
encore vérifier \inserted{d'une part \deleted{les}} les relations
\[
 \left\{\begin{aligned}
 \widetilde h_0r&=\alpha_0,\qquad \underline{p_0\widetilde h_0=\widetilde\beta_0},\\
 \lambda_Xh'_{I'}&=\widetilde h_0s
 \end{aligned}\right.\qquad
 \deleted{\text{(déjà connue, cf }(6.3.10)\text{)}}
 \tag*{(6.3.12)}
\]
% slanted
\margin{NB la première et la troisième de ces relations sont déjà connues, cf (6.3.10)}
établissant la commutativité du
cube supérieur avec le carré diagonal
(diagramme du type (6.3.6), la flèche
pointillée $h_0$ y incluse), d'autre
part les relations
\[
 \left\{\begin{aligned}
 \widetilde i r&=i_0\quad(\text{et }\widetilde i s=\lambda_B\ \text{pour mémoire})\\
 \beta_0\widetilde i&=\widetilde\beta_0,
 \end{aligned}\right.\tag*{(6.3.13)}
\]
où la première ligne est déjà connue
par définition de $\widetilde B$ et de $\widetilde i$. La deuxième
% archive sheet number: 126

\scanpage{128}
\mspage{115}
se ramène aux deux relations
\[
\begin{aligned}
&\underbrace{\begin{array}[t]{c@{\,}c@{\,}c}
 \underline{(\beta_0\widetilde i)r}&=&\underline{\widetilde\beta_0r},\\
 \Vert&&\Vert\\
 \beta_0(\widetilde i r)&&p_0\alpha_0\\
 \Vert&&\\
 \beta_0i_0&&
\end{array}}_{\substack{\text{OK par comm.}\\\text{du carré du bas}\\\text{de (6.3.5)}}}
\quad
\underbrace{\begin{array}[t]{c@{\,}c@{\,}c}
 \underline{(\beta_0\widetilde i)s}&=&\underline{\widetilde\beta_0s}\\
 \Vert&&\Vert\\
 \beta_0(\widetilde i s)&&\lambda_Y\beta_{I'}\\
 \Vert&&\\
 \beta_0\lambda_B&&
\end{array}}_{\substack{\text{OK. par comm.}\\\text{du carré antérieur}\\\text{de (6.3.6).}}}
\end{aligned}
\]
Donc il reste à vérifier \deleted{les} la \inserted{deuxième} \unsure{relations|relation} (6.3.12)
\[
 \deleted{\widetilde h_0 r=\alpha_0,}\quad
 p_0\widetilde h_0=\widetilde\beta_0,\quad
 \deleted{\unsure{\lambda_X h'_{I'}}\ \widetilde h_0 s=\lambda_X h'_{I'}}
\]
ce qui se ramène à
\[
\begin{aligned}
&\begin{array}[t]{c@{\,}c@{\,}c}
 \underline{(p_0\widetilde h_0)r}&=&\widetilde\beta_0 r,\\
 \Vert&&\\
 p_0(\underset{\alpha_0}{\underline{\widetilde h_0r}})&&\\
 \underline{\phantom{p_0\widetilde h_0r}}&&\\[-2pt]
 p_0\alpha_0&&\\
 \multicolumn{3}{c}{\underbrace{\hphantom{(p_0\widetilde h_0)r=\widetilde\beta_0r}}_{\text{OK par (6.3.11)}}}
\end{array}
\qquad
\begin{array}[t]{c@{\,}c@{\,}c}
 \underline{(p_0\widetilde h_0)s}&=&\underline{\widetilde\beta_0s}\\
 p_0(\underset{\lambda_Xh'_{I'}}{\underline{\widetilde h_0s}})&&\lambda_Y\beta_{I'}\\
 \underline{\phantom{(p_0\lambda_X)h'_{I'}}}&&\\[-2pt]
 (\underset{\lambda_Yp_{I'}}{\underline{p_0\lambda_X}})h'_{I'}&&\\
 \underline{\phantom{\lambda_Y(p_{I'}h'_{I'})}}&&\\[-2pt]
 \deleted{\unsure{p_I|\lambda_I}}\unsure{\lambda_Y}(\underset{\unsure{p_{I'}|\beta_{I'}}}{\underline{p_{I'}h'_{I'}}})&&\\
 \multicolumn{3}{c}{\underbrace{\hphantom{(p_0\widetilde h_0)s=\widetilde\beta_0s}}_{\text{OK}}}
\end{array}
\end{aligned}
\]

\deleted{Je dis}
Ceci posé, je dis que les contraintes
(6.3.7), (6.3.8) sur $h_0$ reviennent à la
commutativité des deux triangles dans
le carré

\diagram{\((6.3.14)\)\qquad
\begin{tikzcd}[ampersand replacement=\&]
 \widetilde B\arrow[r,"{\widetilde h_0}"]\arrow[d,"{\widetilde i\in}"',"\Phi"] \& X_0\arrow[d,"{p_0\in\Psi}"]\\
 B_0\arrow[r,"{\beta_0}"]\arrow[ur,dotted,"{h_0}"] \& Y_0,
\end{tikzcd}}
% archive sheet number: 127

\scanpage{129}
\mspage{116}
i.e. aux deux relations
\[
 h_0\widetilde i=\widetilde h_0,\qquad p_0h_0=\beta_0.\tag*{(6.3.15)}
\]
La deuxième de ces relations n'est autre que
la première relation (6.3.7), et il reste à
voir que \deleted{\unsure{6.3.}} la première relation (6.3.15)
équivaut à l'ensemble des deux relations
restantes de (6.3.7), (6.3.8) :
\[
 h_0i_0=\alpha_0,\qquad \lambda_{\deleted{\unsure{I|Y}}X}h'_{I'}=h_0\underset{\widetilde i s}{\underline{\lambda_B}}.\tag*{(6.3.16)}
\]
Or la relation en question $h_0\widetilde i=\widetilde h_0$ équivaut
à l'\unsure{ens} des deux relations
\[
\begin{aligned}
&\underbrace{\begin{array}[t]{c@{\,}c@{\,}c}
 \underline{(h_0\widetilde i)r}&=&\underline{\widetilde h_0 r},\\
 h_0(\underset{i_0}{\underline{\widetilde i r}})&&\alpha_0\\
 \underline{\phantom{h_0(\widetilde i r)}}&&\\[-2pt]
 h_0i_0&&
\end{array}}_{\substack{\text{équivalent à}\\ h_0i_0=\alpha_0}}
\quad
\underbrace{\begin{array}[t]{c@{\,}c@{\,}c}
 \underline{(\overset{\deleted{\unsure{\sim}}}{h}_0\widetilde i)s}&=&\underline{\widetilde h_0 s}\\
 \Vert&&\lambda_Xh'_{I'}\\
 h_0(\underset{\lambda_B}{\underline{\widetilde i s}})&&\\
 \underline{\phantom{h_0(\widetilde i s)}}&&\\[-2pt]
 h_0\lambda_B&&
\end{array}}_{\substack{\text{équivalent à}\\ h_0\lambda_B=\lambda_X \unsure{h''_I|h'_{I'}}}},
\end{aligned}
\]
c'est OK.

Ceci posé, dans le carré (6.3.14), les
hypothèses $i\in\Phi(I)$, $p\in\Psi\langle I\rangle$ impliquent
$\widetilde i_{\deleted{\unsure{0|2}}}\in\Phi$, $p_0\in\Psi$. Donc la flèche cherchée
$h_0$ existe, qed.

Cela achève la démonstration du cor. 6.3,
et par là \unsure{même|aussi}, celle de la dernière
partie, d), de 6.1. On est heureux !
% archive sheet number: 128

\scanpage{130}
\mspage{117}
\uline{Proposition 6.4.} Soit $M$ une catégorie,
et $\Phi\longleftrightarrow\Psi$ un couple $Q$-clos en \deleted{\unsure{la}}
dualité $(\Phi=\Psi^*,\ \Psi=\Phi_*)$. Soit
$u\colon I\longrightarrow I'$ une flèche dans $\mathrm{Cat}$, d'où
suite de foncteurs adjoints

% slanted
\margin{Sous réserve d'existence des $\underrightarrow{\lim}$ et $\underleftarrow{\lim}$ pertinentes dans $M$. \deleted{\unsure{c'est}} pour simplifier $M$ \unsure{stable} pour $\underrightarrow{\lim}$ et $\underleftarrow{\lim}$.}
\diagram{\begin{tikzcd}[ampersand replacement=\&,column sep=huge]
 M(I)\arrow[r,shift left=2.2ex,"{u_!}"]\arrow[r,shift right=2.2ex,"{u_*}"'] \& M(I')\arrow[l,"{u^*}"'].
\end{tikzcd}}

On définit $\Phi(I),\Psi\langle I\rangle,\Phi(I'),\Psi\langle I'\rangle$
comme dans 6.\unsure{2}.5. On a alors :

\textcircled{\scriptsize a}
\[
 \left\{\begin{aligned}
 u_!(\Phi(I))&\subset\Phi(I')\\
 u^*(\unsure{\Psi|\mathscr{T}|T}\langle I'\rangle))&\subset\Psi\langle I'\rangle.
 \end{aligned}\right.\tag*{(6.4.1)}
\]

\textcircled{\scriptsize b} \deleted{D'autre part,} On a l'équivalence
\[
 u^*(\Phi(I'))\subset\Phi(I)\ \Longleftrightarrow\
 u_*(\Psi\langle I\rangle)\subset\Psi\langle I'\rangle.\tag*{(6.4.2)}
\]

\textcircled{\scriptsize c} Les relations équivalentes (6.4.\unsure{2|1}) ont lieu
dans chacun des trois cas suivants :
\deleted{quand $I$ est discrète, ou}

$1^\circ$) $u$ est \unsure{coétale}

$2^\circ$) $I$ est discrète

$3^\circ$) Pour \unsure{tt} $i'\in\operatorname{Ob} I'$, \deleted{\unsure{toute|les}} \unsure{composante|composantes}
\unsure{connexe|connexes} de $I_{/i'}$ admet un objet final.

\uline{Corollaire} 6.5. Soit \deleted{\unsure{le}} $(W,C,F)$ un triple de
Quillen clos \unsure{séparant} dans $M$, et définissons,
pour tt $I$ dans $\mathrm{Cat}$ $C(I)$, $TC(I)$ et $F(I)$, $TF(I)$
comme dans le th. 6.1. a). On a alors
% archive sheet number: 129

\scanpage{131}
\mspage{118}
les équivalences
\[
\bigl(u^*(C(I'))\subset C(I)\bigr)\Longleftrightarrow
\bigl(u_*(TF)\subset TF'\bigr)\unsure{,}
\tag*{(6.5.\unsure{5|1})}
\]
\[
\bigl(u^*(TC(I')\subset TC(I)\bigr)\Longleftrightarrow
\bigl(u_*(F)\subset F'\bigr).
\tag*{(6.5.2)}
\]
Les conditions équivalentes (6.5.1) ainsi que
(6.5.2) sont \unsure{satisfaites} quand on est
dans l'un des cas 1°) 2°) 3°) de
6.4. c).

\uline{Dém de 6.4.} Les assertions \textcircled{a}, \textcircled{b}
sont $\pm$ tautologiques, par adjonction,
comparer $\overline{\underline{\mathrm{XIII}}}$ \unsure{prop} 2 et \unsure{cor}, page 46.
Pour prouver c), il suffit de traiter le
cas 3°, car \deleted{1° et 2°} $1^\circ\Rightarrow3^\circ$, $2^\circ\Rightarrow3^\circ$.
Mais le calcul de \deleted{$u_*(X)$}
\[
u_*(X)(i')\simeq\underset{\unsure{i\in|\mathrm{Ob}} I_{/i'}}{\underleftarrow{\lim}}X(i)
\]
montre que si les composantes
connexes $(I_{/i'})_\alpha$ de $I_{/\unsure{i'|j'}}$ ont
chacune un objet final $e_\alpha$, alors
\[
u_*(X)(i')\simeq\prod_\alpha X(e_\alpha).
\]
Or $\Psi$ \inserted{étant clos à droite} est stable par produits quelconques,
donc $X\in\Psi\langle I\rangle$ implique que
\deleted{\unsure{tous les|tt les}} les $u_*(X)(i')$ sont dans $\Psi$, donc
$u_*(X)$ est dans $\Psi\langle I'\rangle$, \unsure{qed.}
% archive sheet number: 130

\scanpage{132}
$\overline{\underline{\mathrm{XVIII}}}$

\deleted{Modèles (4)}

+ \unsure{Para topos|Paratopos}, catégories accessibles, filtrations cardinales

\uline{NB} Constructions \unsure{transfinies} se trouvent :
\diagram{\[
\begin{array}{lll}
4.6.1\,\ldots & \text{---\quad --\quad ---\quad --} & \text{p.}\;139\text{--}144\\
4.7.\unsure{8|5}.2. & \text{--\quad --\quad --\quad --\quad --} & \text{p.}\;16\unsure{5|8}\text{--}171\\
4.10.3. & \text{--\quad --\quad ---\quad --\quad --} & \text{p.}\;247\text{--}262\\
\inserted{4.13.\unsure{10|20}} & \inserted{\text{--\quad --\quad --\quad --}} &
\inserted{\text{p.}\;295\text{--}302\ \text{bis}}\\
\left\{\begin{array}{l}4.\unsure{13|12}.\unsure{15|18}\\4.13.17\end{array}\right\} &
\text{---\quad --\quad --\quad --} & \text{p.}\;
\deleted{\unsure{305|308}\text{--}\unsure{320|330}}\unsure{308|305}\text{--}330
\end{array}
\]}
% archive sheet number: 131

